\section{Background}

\subsection{Paper 11's Negative Feasibility Result}

Paper~11~\cite{paper11tau} swept the threshold $\tau$ of the Paper~10
adaptive detector~\cite{paper10adaptive} over five values
$\{0.02, 0.03, 0.05, 0.075, 0.10\}$ and found no $\tau$ in the grid
simultaneously satisfies both Paper~10 pre-registered tolerances: at
every $\tau$, the K=50 cell-mean adaptive--lag1 delta was between
$-0.041$ and $-0.053$ pp, missing the Paper~10 $-0.02$~pp ceiling.
The K=200 worst-case was uniformly $0.0$ but the K=50 floor was
too high.

Paper~11's discussion identified three richer detector families as
candidates for reaching the feasibility region: (i) capacity-aware
thresholding, (ii) CUSUM, (iii) Bayesian online change-point. The
present paper evaluates the simplest of these: a single per-K $\tau$
vector locked before evaluation.

\subsection{Why Capacity-Aware First}

Capacity-aware thresholding is the simplest of the three candidate
families. It requires only that the deployment know the operating
capacity --- a quantity already exposed by every patch-management
system. CUSUM and Bayesian detectors require accumulator state and
prior-distribution choices respectively, both of which would expand
the pre-registration surface.

\subsection{Mechanism}

At each window the detector compares the magnitude of the
calibration-target shift $|\Delta_w|$ against $\tau$ and fires if
the threshold is crossed; on fire, the strategy reverts to Paper~4
fixed weights. Paper~11 showed that for a single $\tau$ applied at
all capacities, the firing fraction at K=50 was always too high
because the shifts at that capacity routinely crossed even
$\tau = 0.10$.

A larger $\tau_{50}$ would suppress those false fires; a smaller
$\tau_{200}$ would preserve aggressive firing at high capacity. The
question is whether a single such locked vector can reach the joint
feasibility region.
